\begin{proofbox}
\textbf{\textcolor{CProof}{证明目标与前提}}
设 $C:\dfrac{x^2}{a^2}+\dfrac{y^2}{b^2}=1\ (a>b>0)$，
$P(x_0,y_0)\in C$ 且 $x_0y_0\ne0$。证
\[
     k_{OP}k_l=-\frac{b^2}{a^2}=e^2-1.
\]
下面给出五种角度。\textbf{方法一最适合高考现场；方法四最适合理解几何本质。}

\medskip
\textbf{\textcolor{CProof}{方法一：已知切点式，三行完成（考试最快）}}
\textbf{已知知识：} 椭圆上一点 $P(x_0,y_0)$ 的切线方程是
\[
      \frac{x_0x}{a^2}+\frac{y_0y}{b^2}=1.
\]
\textbf{第 1 步：} 写成斜率式（$y_0\ne0$）：
\[
 y=-\frac{b^2x_0}{a^2y_0}x+\frac{b^2}{y_0},
 \qquad k_l=-\frac{b^2x_0}{a^2y_0}.
\]
\textbf{第 2 步：} 由 $O(0,0)$ 和 $P(x_0,y_0)$ 得
$k_{OP}=\dfrac{y_0}{x_0}$。
\textbf{第 3 步：} 相乘并用 $c^2=a^2-b^2$：
\[
 \boxed{k_{OP}k_l
 =\frac{y_0}{x_0}\left(-\frac{b^2x_0}{a^2y_0}\right)
 =-\frac{b^2}{a^2}=\frac{c^2}{a^2}-1=e^2-1.}
\]
\textbf{考场缩写：} $k_l=-\dfrac{b^2x_0}{a^2y_0}$，
$k_{OP}=\dfrac{y_0}{x_0}$，故积为 $-b^2/a^2=e^2-1$。

\medskip
\textbf{\textcolor{CProof}{方法二：不依赖切点式，隐函数求导（独立验证）}}
将 $y$ 视为 $x$ 的函数，在 $y_0\ne0$ 的附近对椭圆方程求导：
\[
 \frac{2x}{a^2}+\frac{2y}{b^2}y'=0
 \quad\Longrightarrow\quad
 y'=-\frac{b^2x}{a^2y}.
\]
代入 $P(x_0,y_0)$，得
\[
 k_l=-\frac{b^2x_0}{a^2y_0}.
\]
再乘 $k_{OP}=y_0/x_0$，便得
\[
 k_{OP}k_l=-\frac{b^2}{a^2}=e^2-1.
\]
\textbf{为什么可用导数？} 平滑曲线在一点的导数就是该点切线的斜率。
注意求导时 $y$ 是随 $x$ 变化的，所以 $y^2$ 的导数为 $2yy'$。

\medskip
\textbf{\textcolor{CProof}{方法三：参数方程，正切与余切相消}}
令
\[
 x=a\cos\theta,\qquad y=b\sin\theta
 \quad (\sin\theta\cos\theta\ne0).
\]
切点 $P=(a\cos\theta,b\sin\theta)$，故径向斜率
\[
 k_{OP}=\frac{b\sin\theta}{a\cos\theta}
        =\frac ba\tan\theta.
\]
参数曲线切线斜率等于 $(dy/d\theta)/(dx/d\theta)$：
\[
 \frac{dx}{d\theta}=-a\sin\theta,\qquad
 \frac{dy}{d\theta}=b\cos\theta,
\]
\[
 k_l=\frac{b\cos\theta}{-a\sin\theta}
       =-\frac ba\cot\theta.
\]
相乘即
\[
 k_{OP}k_l=-\frac{b^2}{a^2}
 (\tan\theta\cot\theta)=-\frac{b^2}{a^2}.
\]
\textbf{记忆点：} 一个 $\tan\theta$，一个 $-\cot\theta$，
正好互为倒数，只余 $-(b/a)^2$。

\medskip
\textbf{\textcolor{CProof}{方法四：由圆拉伸成椭圆（数形结合、几何本质）}}
\textbf{第 1 步：} 从单位圆 $u^2+v^2=1$ 出发，取
$Q(u_0,v_0)$，且 $u_0v_0\ne0$。圆心到 $Q$ 的半径斜率为
$k_r'=v_0/u_0$；圆的切线垂直半径，所以其斜率为
$k_t'=-u_0/v_0$，从而 $k_r'k_t'=-1$。

\textbf{第 2 步：} 进行拉伸变换
\[
       (u,v)\longmapsto(x,y)=(au,bv).
\]
单位圆变为椭圆，因为 $u=x/a,v=y/b$ 代回得到
$x^2/a^2+y^2/b^2=1$。$Q$ 变成
$P(au_0,bv_0)$。圆的切线也变成椭圆的切线：
事实上，圆的切点式 $u_0u+v_0v=1$ 在代换后变成
$\dfrac{x_0x}{a^2}+\dfrac{y_0y}{b^2}=1$。

\textbf{第 3 步：} 任意非竖直直线在这个变换下，横坐标变化量乘
$a$，纵坐标变化量乘 $b$，因此斜率统一乘 $b/a$。
故
\[
 k_{OP}=\frac ba k_r',\qquad
 k_l=\frac ba k_t',
\]
\[
 \boxed{k_{OP}k_l=\left(\frac ba\right)^2k_r'k_t'
       =-\frac{b^2}{a^2}=e^2-1.}
\]
\textbf{关键直觉：} 不是圆的垂直关系“消失”，而是横纵不等比例
拉伸后，角度不再保持；两个斜率各乘一次 $b/a$。

\medskip
\textbf{\textcolor{CProof}{方法五：法向量与方向向量（联系直线基础知识）}}
设切线方程
\[
      \frac{x_0}{a^2}x+\frac{y_0}{b^2}y=1.
\]
已知 $Ax+By=C$ 的法向量是 $(A,B)$，所以切线法向量为
$\boldsymbol n=(x_0/a^2,y_0/b^2)$。
与 $\boldsymbol n$ 垂直的一个切向量是
\[
 \boldsymbol t=\left(\frac{y_0}{b^2},-\frac{x_0}{a^2}\right),
 \quad
 \boldsymbol n\cdot\boldsymbol t
 =\frac{x_0y_0}{a^2b^2}-\frac{x_0y_0}{a^2b^2}=0.
\]
沿切向量的斜率是纵分量除以横分量，因此
\[
 k_l=\frac{-x_0/a^2}{y_0/b^2}
     =-\frac{b^2x_0}{a^2y_0}.
\]
乘 $k_{OP}=y_0/x_0$ 即得到所证结论。
\textbf{与圆对比：} 当 $a=b$，法向量与 $\overrightarrow{OP}$
同方向；当 $a\ne b$，一般不再同方向。

\medskip
\textbf{\textcolor{CProof}{边界严格性与补充}}
若 $P=(\pm a,0)$，则 $OP$ 水平、切线 $x=\pm a$ 竖直，
$k_l$ 不存在；若 $P=(0,\pm b)$，则 $OP$ 竖直、切线
$y=\pm b$ 水平，$k_{OP}$ 不存在。
\textbf{这些点不能写 $k_{OP}k_l=e^2-1$，更不能写 $0\times\infty=0$。}
但对于趋近这些顶点且不在坐标轴上的点，两斜率\emph{乘积的极限}
仍是 $-b^2/a^2$；顶点自身则用切点式处理。

\medskip
\textbf{\textcolor{CProof}{最终结论}}
\[
 \boxed{\quad x_0y_0\ne0\ \Longrightarrow\
 k_{OP}k_l=-\frac{b^2}{a^2}=e^2-1.\quad}
\]
考试优先“\textbf{切点式 $\to$ 切线斜率 $\to$ 乘积}”；
理解优先“\textbf{圆的切线垂直半径 $\to$ 两方向伸缩}”。
\end{proofbox}
